% Owner: Prateek Shakya
\subsection{Discussion}
Pipelining changes execution time through two factors: cycle count and clock
frequency. For a program with $N$ instructions, execution time is
$C/F_{\max}$, and the speed-up is
\[
\frac{C_{\mathrm{single}}/F_{\mathrm{single}}}
     {C_{\mathrm{pipeline}}/F_{\mathrm{pipeline}}}
= \frac{C_{\mathrm{single}}}{C_{\mathrm{pipeline}}}
  \frac{F_{\mathrm{pipeline}}}{F_{\mathrm{single}}}.
\]
The measured frequency ratio is $56.55/28.43=1.9891$, while the benchmark cycle
ratio ranges from $785/1095=0.7169$ for \texttt{fib} to $786/1238=0.6349$ for
\texttt{gcd}. Consequently, the benchmark speed-ups in
\texttt{analysis/results/speedup.csv} range from 1.2629 for \texttt{gcd} to
1.4260 for \texttt{fib}.

The pipeline also adds a predictable base of four cycles for fill/retirement.
Load-use stalls and taken branches or jumps add $S$ and $2T$, respectively.
This is visible in the measured data: \texttt{bench\_sort} executes 920
instructions in 1444 cycles with 120 stalls and 200 taken transfers, giving
CPI $1444/920=1.570$, while \texttt{bench\_memcpy} has 128 stalls and 189 taken
transfers and reaches CPI $1741/1231=1.414$. The measurements agree with
$C=N+4+S+2T$ for every pipeline run.

The 56.55~MHz pipeline and 28.43~MHz single-cycle result on the ECP5 E85 target
demonstrate a hardware-versus-frequency trade-off. The pipeline uses 3300 LUTs,
1358 FFs and 4 BRAMs; the baseline uses 7392 LUTs, 1040 FFs and 2 BRAMs. The
pipeline is smaller in logic and faster, but registers and forwarding prevent
it from remaining faster on programs with severe branch or load-use overhead.
The tiny \texttt{test4\_loaduse} and \texttt{test5\_branch\_jump} programs do
not show a positive speed-up (0.9945 and 0.9513), because their cycle reductions
are too small to compensate for workload overhead.

These conclusions are limited by the benchmark programs, the supported RV32I
subset and the selected ECP5 target. The benchmarks contain no cache traffic,
operating-system activity or RV32M instructions, and the synthesis count depends
on the OSS CAD tool versions, seed and target configuration. Maximum frequency
is a post-place-and-route timing result, not a silicon measurement. The
single-cycle and pipelined RTL are therefore evaluated as a controlled design
comparison rather than as a claim about general-purpose processor performance.
